\section{Repaso del caso Big Sleep: el ciclo completo, de la hipótesis a la evidencia de fallo}

\subsection{Objetivo del proyecto y definición de la tarea}
Big Sleep (en el curso también se relaciona con Project Naptime) se centra en un objetivo que puede verificarse de forma explícita: dado un repositorio de código, generar una entrada que provoque un fallo detectado por el sanitizer. Este objetivo convierte «descubrir una vulnerabilidad» en un problema de búsqueda ejecutable, lo que permite tanto la evaluación automática como conservar la fidelidad a la práctica de ingeniería real.

\begin{figure}[H]
\centering
\includegraphics[width=0.88\textwidth]{frame-06-big-sleep.jpg}
\caption{Pantalla de presentación del proyecto Big Sleep: convertir el descubrimiento de vulnerabilidades en una tarea de disparo ejecutable \protect\footnotemark}
\end{figure}
\footnotetext{Intervalo de tiempo en el video: 01:14:58--01:15:35.}

\begin{importantbox}{El modelado clave de Big Sleep}
El objetivo no es que el modelo explique de palabra que ``aquí hay una vulnerabilidad'', sino que el sistema produzca una entrada reproducible que haga que el programa genere una anomalía objetiva bajo el sanitizer o el depurador. Dicho de otro modo, la conclusión debe ser reconocida por el entorno de ejecución.
\end{importantbox}

\subsection{Tres categorías centrales de herramientas}
El curso mostró tres categorías de herramientas de este sistema:
\begin{itemize}
    \item Herramienta de exploración de código: búsqueda por símbolo, salto a la definición, seguimiento de referencias;
    \item Herramienta de generación de entradas: que el Agente escriba un programa en Python para construir entradas complejas, en lugar de escribir a mano cadenas de bytes largas;
    \item Herramienta de ejecución con depuración: ejecutar el programa objetivo y devolver los valores de las expresiones correspondientes a watchpoints/breakpoints.
\end{itemize}

\begin{knowledgebox}{Por qué hacer que el Agente genere el programa de entrada en Python}
En escenarios de seguridad, las entradas suelen contener estructuras binarias, formatos anidados o campos de verificación. Que el Agente escriba un ``programa generador'' es más robusto que producir directamente bytes crudos, y también facilita las modificaciones iterativas posteriores.
\end{knowledgebox}

\begin{warningbox}{Tratar el depurador como una ventana de chat interactiva no es eficiente}
El diseño de herramientas mostrado en clase no se detiene en cada paso al estilo de la depuración manual paso a paso, sino que sigue el esquema de ``ejecutar el programa completo y devolver los valores de los puntos de observación alcanzados''. Esto reduce la sobrecarga de interacción y a la vez conserva la evidencia clave, lo que resulta más adecuado para un ciclo automatizado.
\end{warningbox}

\subsection{Resolución de problemas mediante trayectorias}
La trayectoria que muestra el curso es muy representativa: primero se explora el código y se formula una hipótesis de vulnerabilidad, después se produce un informe de la vulnerabilidad junto con una idea de explotación; a continuación se genera la entrada y se ejecuta; si no se produce el fallo, se atribuye la causa del fracaso, se corrige la entrada y se vuelve a ejecutar, hasta provocar la anomalía. Este proceso refleja la ventaja central del Agente: \textbf{el fracaso no es el final, sino la señal de supervisión para la siguiente ronda de búsqueda}.

\begin{figure}[H]
\centering
\includegraphics[width=0.88\textwidth]{frame-08-wrapup.jpg}
\caption{Síntesis de cierre de la clase sobre el estado actual y las perspectivas de los Agentes de seguridad \protect\footnotemark}
\end{figure}
\footnotetext{Intervalo de tiempo en el video: 01:25:45--01:26:20.}

\begin{importantbox}{Un hecho clave a nivel de resultados}
El curso reveló un resultado del mundo real: se descubrió y se divulgó una vulnerabilidad real en SQLite. Más importante aún, el proyecto también mostró capacidad de variant analysis, es decir, buscar variantes similares dentro del mismo proyecto a partir de patrones de vulnerabilidades ya conocidos.
\end{importantbox}

\subsection{Comparación con el fuzzing tradicional}
El expositor mencionó un detalle muy convincente: incluso conociendo de antemano la ubicación de la vulnerabilidad, al intentar encontrarla en sentido inverso con un fuzzer genérico, una ejecución prolongada (el curso menciona un orden de magnitud de unas 150 horas) sigue teniendo dificultades para dar directamente con la misma vulnerabilidad. Esta comparación sugiere que la guía semántica del Agente puede aportar una ganancia en la exploración de rutas profundas.

\begin{knowledgebox}{El valor potencial de la búsqueda agentic}
\begin{itemize}
    \item Puede aprovechar conocimiento previo (vulnerabilidades históricas, semántica de los parches) para realizar una búsqueda dirigida;
    \item Puede reescribir la estrategia de generación de entradas después de un fracaso;
    \item Puede combinarse con la retroalimentación de la depuración para reducir el espacio de rutas sospechosas.
\end{itemize}
\end{knowledgebox}

\begin{warningbox}{No hay que extrapolar el éxito de un solo caso a una ``superioridad total''}
Big Sleep demostró la viabilidad, pero eso no significa que todos los proyectos y todos los tipos de vulnerabilidades ya puedan automatizarse. La conclusión de la clase es que ``apenas está empezando y hay mucho margen'', no que ``el problema ya esté resuelto''.
\end{warningbox}

\subsection{Resumen de la sección}
El valor de Big Sleep está en que ejemplifica un ciclo completo: el objetivo es ejecutable, las herramientas son invocables, la trayectoria es iterable y el resultado es verificable. Hace avanzar a los Agentes de seguridad, de la prueba de concepto a una práctica de ingeniería reproducible.
